\subsection{正规化引理}

在本节及以下各节中，k 表示一个交换域。

定理 1（正规化引理）. 设 A 为一个有限生成 \(k\)-代数，并设 \(\mathfrak{a}_1 \subset \mathfrak{a}_2 \subset \ldots \subset \mathfrak{a}_p\) 为 A 的理想的递增有限序列，满足 \(p \geqslant 1\) 且 \(\mathfrak{a}_p \neq \mathbf{A}\)。那么存在 A 中元素的一个有限序列 \((x_i)_{1 \leqslant i \leqslant n}\)，它们在 \(k\) 上代数无关（第 III 章 § 1 第 1 节），并且满足：

\begin{enumerate}
\def\labelenumi{(\alph{enumi})}
\item
  A 在环 \(\mathbf{B} = k[x_1, \ldots, x_n]\) 上是整的。
\item
  对所有 \(j\)（满足 \(1 \leqslant j \leqslant p\)），存在整数的一个递增序列 \((h(j))_{1 \leqslant j \leqslant p}\)，使得对所有 \(j\)，理想 \(\mathfrak{a}_j \cap \mathbf{B}\)（\(\mathbf{B}\) 中的）由 \(x_1, \ldots, x_{h(j)}\) 生成。
\end{enumerate}

首先注意，只需在 A 是多项式代数 \(k[Y_{1},\ldots,Y_{m}]\) 的情形证明本定理即可。因为在一般情形，A 同构于这样一个代数 \(A'\) 除以理想 \(a_{0}'\) 所得的商；用 \(a_{j}'\) 记 \(a_{j}\) 在 \(A'\) 中的原像，并设 \(x_{i}'(1\leqslant i\leqslant r)\) 是 \(A'\) 的元素，它们对环 \(A'\) 与理想的递增序列 \(a_{0}'\subset a_{1}'\subset\cdots\subset a_{p}'\) 满足命题的条件。于是，像 \(x_{i}\)（即 \(x_{i}'\) 在 A 中的像，其中 \(i>h(0)\)）满足所要的条件；这对条件 (b) 是显然的，而对条件 (a) 则由 §1 第 1 节命题 2 得出；最后，若诸 \(x_{i}(h(0)+1\leqslant i\leqslant r)\) 在 k 上不是代数无关的，就会存在一个非零多项式

\[
\mathrm{Q} \in k [ \mathrm{X} _ {h (0) + 1}, \dots , \mathrm{X} _ {r} ]
\]

使得 \(\mathrm{Q}(x_{h(0)+1}^{\prime},\ldots,x_{r}^{\prime})\in\mathfrak{a}_{0}^{\prime}\cap\mathbf{B}^{\prime}\)，其中记 \(B^{\prime}=k[x_{1}^{\prime},\ldots,x_{r}^{\prime}]\)；但按假设，\(a_{0}^{\prime}\cap B^{\prime}\) 的每个元素都可以唯一地写成系数在 k 中的、诸 \(x_{j}^{\prime}\ (1\leqslant j\leqslant r)\) 的多项式，其每个单项式都至少含有某个 \(x_{j}^{\prime}\)（其中 \(1\leqslant j\leqslant h(0)\)）；于是得到矛盾，这就证明了我们的断言。

因此，在证明的余下部分我们将设 \(A = k[Y_{1}, \ldots, Y_{m}]\)，并对 p 作归纳论证。

\begin{enumerate}
\def\labelenumi{(\Alph{enumi})}
\tightlist
\item
  p = 1。我们区分两种情形：
\end{enumerate}

\textbf{(A1) 理想 \(\mathfrak{a}_1\) 是由一个元素 \(x_1 \notin k\) 生成的主理想。}\par

\(x_{1} = \mathrm{P}(\mathrm{Y}_{1},\ldots ,\mathrm{Y}_{m})\)，其中 \(\mathbf{P}\) 是非常数多项式。我们将看到，适当选取整数 \(r_i > 0\)，环 A 在 \(\mathbf{B} = k[x_1,x_2,\dots ,x_m]\) 上是整的，其中

\[
x _ {i} = \mathrm{Y} _ {i} - \mathrm{Y} _ {1} ^ {r _ {i}} \quad (2 \leqslant i \leqslant m).\tag{1}
\]

为此只需选取诸 \(r_i\) 使 \(Y_1\) 在 \(B\) 上是整的（§ 1 第 1 节命题 4）。现在，有关系式

\[
\mathrm{P} \left(\mathrm{Y} _ {1}, x _ {2} + \mathrm{Y} _ {1} ^ {r _ {2}}, \dots , x _ {m} + \mathrm{Y} _ {1} ^ {r _ {m}}\right) - x _ {1} = 0.\tag{2}
\]

设 \(\mathbf{P} = \sum_{\mathbf{p}} a_{\mathbf{p}} \mathbf{Y}^{\mathbf{p}}\)，其中 \(\mathbf{p} = (p_1, \ldots, p_m)\)，诸 \(p_i\) 是整数 \(\geqslant 0\)，\(\mathbf{Y}^{\mathbf{p}} = \mathbf{Y}_1^{p_1} \ldots \mathbf{Y}_m^{p_m}\)，诸 \(a_{\mathbf{p}}\) 皆 \(\neq 0\) 且属于 \(k\)，并且至少有一个指标 \(\mathbf{p}\) 异于 \((0, \ldots, 0)\)；关系式 (2) 可以写成

\[
\sum_ {\mathbf {p}} a _ {\mathbf {p}} \mathrm{Y} _ {1} ^ {p _ {1}} (x _ {2} + \mathrm{Y} _ {1} ^ {r _ {2}}) ^ {p _ {2}} \dots (x _ {m} + \mathrm{Y} _ {1} ^ {r _ {m}}) ^ {p _ {m}} - x _ {1} = 0.\tag{3}
\]

记 \(f(\mathbf{p}) = p_{1} + r_{2}p_{2} + \cdots + r_{m}p_{m}\)，并设诸 \(r_{i}\) 选取成使诸 \(f(\mathbf{p})\) 互不相同（例如取 \(r_{i} = h^{i}\) 即可，其中 h 是一个严格大于所有 \(p_{j}\) 的整数（《集合论》第 III 章 §5 第 7 节命题 8））。那么将存在唯一的组 \(\mathbf{p} = (p_{1}, \ldots, p_{m})\) 使 \(f(\mathbf{p})\) 达到最大，且关系式 (3) 可以写成

\[
a _ {\mathbf {p}} \mathrm{Y} _ {1} ^ {f (\mathbf {p})} + \sum_ {j <   f (\mathbf {p})} \mathrm{Q} _ {j} (x _ {1}, \dots , x _ {m}) \mathrm{Y} _ {1} ^ {j} = 0\tag{4}
\]

其中诸 \(Q_{j}\) 是 \(k[Y_{1},\ldots,Y_{m}]\) 中的多项式；由于 \(a_{p} \neq 0\) 在 k 中可逆，(4) 当然是系数在 B 中的一个整依赖方程，由此即得我们的断言。

于是分式域 \(k(\mathrm{Y}_1, \ldots, \mathrm{Y}_m)\)（即 \(\mathbf{A}\) 的分式域）在分式域 \(k(x_1, \ldots, x_m)\)（即 \(\mathbf{B}\) 的分式域）上是代数的，这证明（《代数》第 V 章 § 5 第 3 节定理 4）诸 \(x_i\)（\(1 \leqslant i \leqslant m\)）是代数无关的。此外，\(a_1 \cap B = Bx_1\)；因为每个元素 \(z \in a_1 \cap B\) 都可写成 \(z = x_1z'\)，其中 \(z' \in A \cap k(x_1, \ldots, x_m)\)；而由于 B 是整闭的，有 \(A \cap k(x_1, \ldots, x_m) = k[x_1, \ldots, x_m] = B\)（§ 1 第 3 节命题 13 的推论 2）；因此 \(z' \in B\)，这就完成了此情形中性质 (a) 与 (b) 的证明。

\textbf{(A2) 一般情形 \((p = 1)\)。}\par

我们对 \(m\) 作归纳论证，\(m = 0\) 的情形是平凡的。显然可设 \(a_1 \neq 0\)（否则可取 \(x_i = Y_i\)（对 \(1 \leqslant i \leqslant m\)）且取 \(h(1) = 0\)）。设 \(x_1\) 为 \(a_1\) 的一个非零元素；由 (A1)，存在 \(t_2, \ldots, t_m\) 使 \(x_1, t_2, \ldots, t_m\) 在 \(k\) 上代数无关，A 在 \(C = k[x_1, t_2, \ldots, t_m]\) 上是整的，且 \(x_1A \cap C = x_1C\)。由归纳假设，存在元素 \(x_2, \ldots, x_m\)（属于 \(k[t_2, \ldots, t_m]\)）与一个整数 \(h\)，使得 \(k[t_2, \ldots, t_m]\) 在 \(B' = k[x_2, \ldots, x_m]\) 上是整的，\(x_2, \ldots, x_m\) 在 \(k\) 上代数无关，且理想 \(a_1 \cap B'\) 由 \(x_2, \ldots, x_h\) 生成。于是 C 在 \(B = k[x_1, x_2, \ldots, x_m]\) 上是整的（\(\S\) 1 第 1 节命题 5 的推论），从而 A 亦然（\(\S\) 1 第 1 节命题 6）；与情形 (A1) 中同样的论证表明 \(x_1, \ldots, x_m\) 在 \(k\) 上代数无关；最后，由于 \(x_1 \in a_1\) 且 \(B = B'[x_1]\)，有 \(a_1 \cap B = Bx_1 + (a_1 \cap B')\)，又由于 \(a_1 \cap B'\)（在 \(B'\) 中）由 \(x_2, \ldots, x_h\) 生成，故 \(a_1 \cap B\)（在 B 中）由 \(x_1, x_2, \ldots, x_h\) 生成。

\textbf{(B) 从 p - 1 到 p 的过渡。}\par

设 \(t_1, \ldots, t_m\) 为 A 的元素，它们对理想的递增序列 \(a_1 \subset \ldots \subset a_{p-1}\) 满足定理的条件，并记 \(r = h(p - 1)\)。由 (A2)，存在元素 \(x_{r+1}, \ldots, x_m\)（属于 \(k[t_{r+1}, \ldots, t_m]\)）与一个整数 \(s\)，使得

\[
\mathbf {C} = k [ t _ {r + 1}, \dots , t _ {m} ]
\]

在 \(\mathbf{B}' = k[x_{r+1},\ldots,x_m]\) 上是整的，\(x_{r+1},\ldots,x_m\) 在 \(k\) 上代数无关，且理想 \(\mathfrak{a}_p\cap\mathbf{B}'\) 由 \(x_{r+1},\ldots,x_s\) 生成。记 \(x_i=t_i\)（对 \(i\leqslant r\)），则所得到的族 \((x_i)_{1\leqslant i\leqslant m}\) 以 \(h(p)=s\) 解决了问题。因为 A 在 \(\mathbf{C}[t_1,\ldots,t_r]=\mathbf{C}[x_1,\ldots,x_r]\) 上是整的，从而 A 在 \(\mathbf{B}=k[x_1,\ldots,x_m]=\mathbf{B}'[x_1,\ldots,x_r]\) 上也是整的，因为 C 在 \(\mathbf{B}'\) 上是整的（§ 1 第 1 节命题 5 的推论与命题 6）；可以像情形 (A1) 中那样证明诸 \(x_i\) 在 \(k\) 上代数无关。另一方面，对 \(j\leqslant p-1\)，理想

\[
a _ {j} \cap k [ x _ {1}, \dots , x _ {r}, t _ {r + 1}, \dots , t _ {m} ]
\]

按假设是 \(x_1, \ldots, x_r, t_{r+1}, \ldots, t_m\) 中的多项式所成之集，这些多项式的每个单项式都含有元素 \(x_1, \ldots, x_{h(j)}\) 之一；由于 \(x_{r+1}, \ldots, x_m\) 是 \(t_{r+1}, \ldots, t_m\) 的系数在 \(k\) 中的多项式，立即可见，\(x_1, \ldots, x_r, x_{r+1}, \ldots, x_m\)（系数在 \(k\) 中）的多项式属于 \(a_j\) 仅当其所有单项式都含有元素 \(x_1, \ldots, x_{h(j)}\) 之一。最后，由于 \(x_1, \ldots, x_r\) 属于 \(a_{p-1}\) 从而也属于 \(a_p, a_p \cap B'[x_1, \ldots, x_r]\)，它由 \(x_1, \ldots, x_r\) 的、系数在 \(B'\) 中且常数项属于 \(a_p \cap B'\) 的多项式组成；因此这个理想由 \(x_1, \ldots, x_r, x_{r+1}, \ldots, x_t\) 生成。

推论 1. 设 A 为整环，B 为一个包含 A 作为子环的有限生成 A-代数。那么存在 A 的元素 \(s \neq 0\) 与 B 的一个子代数 B'，它同构于多项式代数 \(\mathrm{A}[\mathrm{Y}_{1}, \ldots, \mathrm{Y}_{n}]\)，使得 \(\mathrm{B}[s^{-1}]\)（第 II 章 §2 第 1 节）在 \(\mathrm{B}'[s^{-1}]\) 上是整的。

记 \(S = A - \{0\}\)，并令 \(k = S^{-1}A\)，即 \(A\) 的分式域；显然 \(S^{-1}B\) 是一个有限生成 \(k\)-代数，且由于按假设它包含 \(k\)（第 II 章 § 2 第 4 节定理 1），它不退化为 0。于是由定理 1（应用于 \(p = 1\) 与 \(a_1 = 0\)）存在一个有限序列 \((x_i)_{1 \leqslant i \leqslant n}\)（由 \(S^{-1}B\) 的元素组成），它们在 \(k\) 上代数无关，且使 \(S^{-1}B\) 在 \(k[x_1, \ldots, x_n]\) 上是整的。设 \((z_j)_{1 \leqslant j \leqslant m}\) 为 A-代数 \(B\) 的一个生成系；在 \(S^{-1}B\) 中，每个 \(z_j/1\) 都满足一个整依赖方程

\[
(z _ {j} / 1) ^ {q _ {j}} + \sum_ {h <   q _ {j}} \mathrm{P} _ {h j} (x _ {1}, \dots , x _ {n}) (z _ {j} / 1) ^ {h} = 0\tag{5}
\]

其中诸 \(P_{hj}\) 是 \(x_{i}\) 的系数在 k 中的多项式。存在 A 的元素 \(s \neq 0\)，使我们可以写 \(x_{i} = y_{i}/s\)，其中 \(y_{i} \in B\)（对 \(1 \leqslant i \leqslant n\)），并且诸 \(P_{hj}\) 的所有系数都具有 c/s 的形式，其中 \(c \in A\)；最后，如有必要把 s 换成 S 中元素的一个乘积，我们可假定在 B 中

\[
s z _ {j} ^ {q _ {j}} + \sum_ {h <   q _ {j}} Q _ {h j} z _ {j} ^ {h} = 0\tag{6}
\]

其中诸 \(Q_{hj}\) 是 \(y_{1},\ldots,y_{n}\) 的系数在 A 中的多项式；若记 \(z_{j}^{\prime}=sz_{j}\)，则将 (6) 乘以 \(s^{q_{j}-1}\) 即见 \(z_{j}^{\prime}\) 在 \(B^{\prime}=A[y_{1},\ldots,y_{n}]\) 上是整的。我们证明诸 \(y_{i}\) 在 A 上代数无关；若有一个形如 \(\sum_{p}a_{p}y_{1}^{p_{1}}\ldots y_{n}^{p_{n}}=0\) 的关系，其中对所有 p 有 \(a_{p}\in A\)，则可推出 \(\sum_{p}a_{p}^{\prime}x_{1}^{p_{1}}\ldots x_{n}^{p_{n}}=0\)（在 \(S^{-1}B\) 中），其中 \(a_{p}^{\prime}=a_{p}s^{p_{1}+\cdots+p_{n}}\)（在 k 中）；于是按假设对所有 p 有 \(a_{p}^{\prime}=0\)，从而对所有 p 有 \(a_{p}=0\)。此外，在环 \(B[s^{-1}]\) 中，每个 \(z_{j}^{\prime}/1\) 都在 \(B^{\prime}[s^{-1}]\) 上是整的（§1 第 1 节命题 2），且由于 \(z_{j}/1=(z_{j}^{\prime}/1)(1/s)\)（在 \(B[s^{-1}]\) 中），可见诸 \(z_{j}/1\) 都在 \(B^{\prime}[s^{-1}]\) 上是整的，这就完成了证明（§1 第 1 节命题 4）。

推论 2. 设 K 为域，A 为 K 的子环，L 为 A 的分式域。若 K 是有限生成 A-代数，则 {[}K:L{]} 有限，且存在 A 中的 \(a \neq 0\) 使 \(\mathbf{L} = \mathbf{A}[a^{-1}]\)。

由推论 1，存在元素 \(x_{1}, \ldots, x_{n}\)（属于 \(\mathbf{K}\)）与一个元素 \(a \neq 0\)（属于 \(\mathbf{A}\)），使 \(x_{1}, \ldots, x_{n}\) 在 A 上（从而在 L 上）代数无关，且 K 在子环 \(A[x_{1},\ldots,x_{n},a^{-1}]\) 上是整的。于是由 §2 第 1 节引理 2 可知 \(A[x_{1},\ldots,x_{n},a^{-1}]\) 是一个域。但整环 C 上的多项式环 \(C[Y_{1},\ldots,Y_{n}]\) 中仅有的可逆元就是 C 的可逆元；把这个注记应用于 \(C = A[a^{-1}]\)，便知必有 n = 0，且 \(A[a^{-1}]\) 是一个域，按 L 的定义它等于 L。由于 K 在 L 上是整的且是有限生成 L-代数，次数 {[}K:L{]} 有限（§1 第 1 节命题 4）。

推论 3. 设 A 为整环，B 为有限生成 A-代数，b 为 B 中元素，使得 \(zb^n \neq 0\)（对 A 中所有 \(z \neq 0\) 及每个整数 \(n > 0\)）。设 \(\rho: A \to B\) 为典范同态；存在 A 中的 \(a \neq 0\)，使得对每一个同态 \(f\)（从 A 到代数闭域 L），若满足 \(f(a) \neq 0\)，则存在从 B 到 L 的同态 \(g\)，使 \(g(b) \neq 0\) 且 \(f = g \circ \rho\)。

对 \(b\) 所作的假设蕴含：若 \(h\) 是典范同态 \(x \mapsto x / 1\)（从 \(\mathbf{B}\) 到 \(\mathbf{B}[b^{-1}]\)），则同态 \(h \circ \rho\)（从 \(\mathbf{A}\) 到 \(\mathbf{B}[b^{-1}]\)）是单射。于是由推论 1，存在元素 \(a \neq 0\)（属于 \(\mathbf{A}\)）与子环 \(\mathbf{B}'\)（\(\mathbf{B}[b^{-1}]\) 的），使得 \(\mathbf{B}[b^{-1}, a^{-1}]\) 在 \(\mathbf{B}'[a^{-1}]\) 上是整的，且 \(\mathbf{B}'\) 同构于多项式代数 \(\mathbf{A}[\mathbf{Y}_1, \ldots, \mathbf{Y}_n]\)。设 \(f\) 是 \(\mathbf{A}\) 到代数闭域 \(\mathbf{L}\) 的一个满足 \(f(a) \neq 0\) 的同态；存在从 \(\mathbf{A}[\mathbf{Y}_1, \ldots, \mathbf{Y}_n]\) 到 \(\mathbf{L}\) 的延拓 \(f\) 的同态，从而存在一个同态 \(f'\)（从 \(\mathbf{B}'\) 到 \(\mathbf{L}\)），它延拓 \(f\)。由于 \(f'(a) \neq 0\)（在 \(\mathbf{L}\) 中），存在同态 \(f''\)（从 \(\mathbf{B}'[a^{-1}]\) 到 \(\mathbf{L}\)），使得

\[
f ^ {\prime \prime} (x / a ^ {n}) = f ^ {\prime} (x). (f (a)) ^ {- n}
\]

对所有 \(x \in \mathbf{B}'\) 与所有 \(n > 0\)（第 II 章 §2 第 1 节命题 1）。最后，由于 \(\mathbf{B}[b^{-1}, a^{-1}]\) 在 \(\mathbf{B}'[a^{-1}]\) 上是整的，存在同态 \(f'''\)（从 \(\mathbf{B}[b^{-1}, a^{-1}]\) 到 \(\mathbf{L}\)），它延拓 \(f''\)（§2 第 1 节定理 1 的推论 4）。若 \(j: x \mapsto x / 1\) 是 \(\mathbf{B}\) 到 \(\mathbf{B}[b^{-1}, a^{-1}]\) 的典范同态，则 \(g = f''' \circ j\) 解决了问题，因为 \(j(b)\) 在 \(\mathbf{B}[b^{-1}, a^{-1}]\) 中可逆，从而 \(f'''(j(b)) \neq 0\)（在 \(\mathbf{L}\) 中）。

注意，若在推论 3 中设 B 为整环且 A ⊆ B，则对 b 的假设等价于 \(b \neq 0\)。

